\section{Diode}

A diode is a two-terminal, passive (it can not amplify), non-linear
and time-invariant device that allows current to flow in one direction while blocking
it in the opposite direction.

It is constructed from a semiconductor where one side (N) is doped
with donors that give free electrons and the other side (P) is doped
with acceptors that give free holes. Where the two sides meet, carriers
diffuse across and recombine, leaving a depletion region with no free
carriers and a built-in potential,
\[
  V_{bi} = V_T \ln\left(\frac{N_A N_D}{n_i^2}\right)
\]
where \(N_A\), \(N_D\) are the acceptor and donor concentrations and
\(n_i\) is the intrinsic carrier concentration. A forward voltage lowers
this barrier and current grows exponentially; appreciable current flows
once \(V_D\) approaches the knee (threshold) voltage, typically around
0.6--0.7\,V for silicon and 0.2--0.3\,V for germanium diodes.

\begin{enumerate}
  \item \textbf{Types of Diodes}:
    There are many types of diodes, some of the most common are:
    \begin{itemize}
      \item \textbf{Rectifier Diode}: Used to convert AC to DC.\ They
        are very useful in power supply circuits. Examples include 1N4001,
        1N4007, 1N5408, etc.
      \item \textbf{Zener Diode}: Used to regulate voltage. They exploit
        the reverse breakdown region of the diode to maintain a
        constant voltage across the diode. Examples include the
        1N4733A (5.1\,V), 1N4742A (12\,V), and the BZX79 series.
      \item \textbf{Light Emitting Diode (LED)}: Used to emit
        light when current flows through it. Under the photovoltaic
        effect, they can also be used to convert light into electrical
        current. Although they aren't as efficient as photodiodes.
        Examples include the 5\,mm red T-1 3/4 (e.g.\ Kingbright L-53LID),
        and high-power types such as the Cree XLamp XP-G3 and Lumileds
        Luxeon Rebel.
      \item \textbf{Schottky Diode}: Used for fast switching
        applications. It is a metal-semiconductor junction that conducts
        with majority carriers only, so there is no stored minority
        charge to remove (no reverse recovery), which lets it switch
        faster. It also has a lower forward voltage drop (0.15--0.45\,V)
        but a higher reverse leakage than a PN diode. Examples include the 1N5817, 1N5819, and the
        SS14 (surface-mount).
      \item \textbf{Photodiode}: Used to convert light into
        electrical current. Examples include the BPW34, BPW21, and
        the SFH 203 series.
      \item \textbf{Varactor Diode}: Used to vary capacitance in
        a circuit. The varactor diode is a reverse-biased diode that
        has a variable capacitance that changes with the applied voltage.
        Examples include the BB135 and the 1SV149, commonly used in
        RF tuning circuits and voltage-controlled oscillators (VCOs).
      \item \textbf{TVS Diode}: Used to protect circuits from voltage
        spikes. A TVS diode (unidirectional or bidirectional) can clamp
        voltage spikes to a safe level. These are often used to protect
        sensitive electronics from electrostatic discharge (ESD) and
        other transient voltage events. Examples include the P6KE
        series (unidirectional/bidirectional, 600\,W) and the
        SMBJ series (surface-mount).
    \end{itemize}
  \item \textbf{Diode Parameters}:
    The parameters of a diode are:
    \begin{itemize}
      \item \textbf{Forward Voltage Drop (\(V_F\))}: The voltage drop
        across the diode when it is forward biased.
      \item \textbf{Reverse Breakdown Voltage (\(V_{BR}\))}: The voltage
        at which the diode will conduct in reverse bias.
      \item \textbf{Reverse Saturation Current (\(I_S\))}: The scale
        current of the Shockley equation; the reverse current tends to
        \(-I_S\) for \(V_D \ll -nV_T\) (before breakdown).
      \item \textbf{Maximum Forward Current (\(I_F\))}: The maximum
        current that can flow through the diode when it is forward biased.
      \item \textbf{Reverse Leakage Current (\(I_R\))}: The current
        that flows at the rated reverse voltage (in practice much larger
        than \(I_S\) due to leakage and generation in the depletion region).
    \end{itemize}
  \item \textbf{Mathematical Models}:
    There are several mathematical models that can be used to describe
    the behaviour of a diode. The most common models are:
    \begin{itemize}
      \item \textbf{Ideal Model}: The ideal diode model assumes
        that the diode has zero forward voltage drop and infinite reverse
        resistance. It is a good approximation when \(V_F\) is negligible
        against the other voltages of the circuit. It is given by:
        \[
          \begin{cases}
            I = 0 & \text{if } V_D < 0 \hspace{1cm} \text{(open)} \\
            V_D = 0 & \text{if } I > 0 \hspace{1cm} \text{(short)}
          \end{cases}
        \]
      \item \textbf{Non-linear Model}: The Shockley diode model
        is a more accurate model that takes into account the forward voltage
        drop and the reverse saturation current. It is a good approximation
        at low and moderate currents; at high currents the series
        resistance \(R_S\) and high level injection take over.
        \[
          I = I_S \left( e^{\frac{V_D}{n V_T}} - 1 \right)
        \]
        where:
        \begin{itemize}
          \item \(I\) is the current through the diode.
          \item \(I_S\) is the reverse saturation current.
          \item \(V_D\) is the voltage across the diode.
          \item \(n\) is the ideality factor (typically between 1
            and 2).
          \item \(V_T\) is the thermal voltage, given by \(V_T = \frac{kT}{q}\).
        \end{itemize}
        This model does not account for the reverse breakdown conduction of
        the diode, nor does it account for the series or breakdown resistance
        of the diode.
      \item \textbf{Linear Model}: The piecewise linear model approximates
        the current-voltage relationship as a series of linear segments.
        The model also includes the effects of series resistance \(R_S\)
        and breakdown resistance \(R_{ZK}\).
        \[
          I \approx
          \begin{cases}
            \dfrac{V_D + V_{BR}}{R_{ZK}} & \text{if } V_D < -V_{BR} \\[2mm]
            0 & \text{if } -V_{BR} \leq V_D < V_{TH} \\[2mm]
            \dfrac{V_D - V_{TH}}{R_S} & \text{if } V_D \geq V_{TH}
          \end{cases}
        \]
        It is obtained by solving the Non-Linear Model in two separate
        voltage regimes: forward, cutoff. The breakdown regime is added as
        an independent term to account for reverse breakdown conduction.
        \begin{enumerate}
          \item \textbf{Series Resistance \(R_S\):}
            The Non-Linear Model treats the entire applied voltage \(V_D\)
            as appearing across the ideal exponential junction. However In a
            real diode, part of \(V_D\) drops across the bulk and contact
            series resistance \(R_S\), so only the remainder, the junction
            voltage \(V_j = V_D - I R_S\), drives the exponential term:
            \[
              I = I_S \left( e^{\frac{V_D - I R_S}{n V_T}} - 1 \right)
            \]
          \item \textbf{Forward Regime:}
            For \(I \gg I_S\), the \(-1\) term is negligible and
            \(I \approx I_S e^{V_j / (nV_T)}\), so
            \[
              V_j \approx n V_T \ln\!\left(\frac{I}{I_S}\right)
            \]
            Because the logarithm varies slowly, \(V_j\) changes by only
            \(n V_T \ln(10) \approx 60\,\text{mV}\) per decade of current, so
            across a diode's typical operating range \(V_j\) is well
            approximated as constant. Fixing \(V_j \approx V_{TH}\), the
            junction voltage at some nominal operating current \(I_0\), and
            substituting into \(V_j = V_D - I R_S\) gives
            \[
              I \approx \frac{V_D - V_{TH}}{R_S} \hspace{1cm} \text{if } V_D \geq V_{TH}
            \]
          \item \textbf{Cutoff Regime:}
            For \(-V_{BR} \leq V_D < V_{TH}\), the current is small enough that
            \(I R_S\) is negligible, so \(V_j \approx V_D\). Two cases:
            \begin{align*}
              V_D \ll -nV_T &: \hspace{0.5cm} I = I_S \left( e^{V_D / (nV_T)} - 1 \right) \approx -I_S \approx 0 \\
              -nV_T \lesssim V_D < V_{TH} &: \hspace{0.5cm} I = I_S \left( e^{V_D / (nV_T)} - 1 \right) \ll I_S e^{V_{TH}/(nV_T)} \approx I_0 \Rightarrow I \approx 0
            \end{align*}
            since \(I_S\) is typically \(10^{-9}\) to \(10^{-15}\,\text{A}\) and
            every \(60\,\text{mV}\) below \(V_{TH}\) divides the current by 10.
          \item \textbf{Breakdown Regime:}
            Reverse breakdown arises from avalanche multiplication or Zener
            tunneling, mechanisms entirely absent from the minority-carrier
            diffusion current that the Shockley equation models. No asymptotic
            limit of the exponential term can reproduce this behavior, so it
            must be introduced as an independent term. Modeled with the same
            linear structure as the forward regime, using breakdown resistance
            \(R_{ZK}\):
            \[
              I \approx \frac{V_D + V_{BR}}{R_{ZK}} \hspace{1cm} \text{if } V_D < -V_{BR}
            \]
        \end{enumerate}
          \item \textbf{Small Signal Model}: Linearizing the Shockley model
        about \((V_{DQ}, I_{DQ})\) (Section~\ref{sec:analysis_linearization}):
        \begin{align*}
          g_d &= \left.\frac{dI}{dV_D}\right|_{Q} \\
          g_d &= \frac{d}{dV_D}\left[I_S\left(e^{\frac{V_D}{nV_T}} - 1\right)\right]_{Q} \\
          g_d &= \frac{I_S e^{\frac{V_{DQ}}{nV_T}}}{nV_T} \\
          g_d &= \frac{I_{DQ} + I_S}{nV_T} \approx \frac{I_{DQ}}{nV_T}
        \end{align*}
        In parallel with \(g_d\) sit the junction (depletion) capacitance
        and the diffusion capacitance,
        \begin{align*}
          C_j &= \frac{C_{j0}}{\left(1 - \frac{V_{DQ}}{V_{bi}}\right)^{m}} \hspace{1cm} [m = 0.5 \text{ abrupt},\ m = 0.33 \text{ graded}] \\
          C_d &= T_F\, g_d \hspace{1cm} [T_F = \text{forward transit time}]
        \end{align*}
        giving the Q4 admittance \(Y_d(s) = g_d + s(C_j + C_d)\). A reverse
        biased diode has \(g_d \approx 0\) and \(C_d \approx 0\), leaving
        \(Y_d = sC_j\), which is the varactor.
    \end{itemize}
  \item \textbf{Circuit Analysis:}
    In Q5 use the Non-Linear (Shockley) or the piecewise Linear model to
    find the bias point. In Q4 replace the diode by \(Y_d(s)\) evaluated
    at that bias point.
\end{enumerate}
